\Subhead{3.10.1}{Генерация второй гармоники, самофокусировка и самомодуляция частично когерентных волн. Преобразование спектров шумовых волн в нелинейных средах без дисперсии}

Нелинейный отклик определяется высшими моментами поля. При локальной генерации второй гармоники без истощения $E_2\propto E_1^2$, следовательно
\begin{equation}
 \langle I_2\rangle\propto\langle I_1^2\rangle.
\end{equation}
Для постоянной когерентной амплитуды $\langle I^2\rangle=\langle I\rangle^2$, для нулевого кругового гауссовского поля $\langle I^2\rangle=2\langle I\rangle^2$. Фактор два относится к выбранному локальному, неистощаемому и согласованному режиму. Конечная длина среды, фазовый синхронизм, полоса и времена отклика могут изменить сравнение. Взаимная когерентность второй гармоники в той же модели $\Gamma_2(1,2)\propto2\Gamma_1^2(1,2)$, её спектр — свёртка спектра основного поля.

При Керровском отклике $\Delta\phi=k_0n_2I L$: флуктуации интенсивности вызывают фазовую самомодуляцию и изменение когерентности. Подстановка $I\to\langle I\rangle$ не замыкает уравнение среднего поля: $\langle|E|^2E\rangle$ не определяется одним $\langle E\rangle$. Нужна статистика поля и сравнение времени отклика среды со временем когерентности. Медленный отклик усредняет реализации; быстрый реагирует на спекловые пики, включая их локальную самофокусировку.

Простейший локальный пример изменения спектра даёт безмемориальная характеристика. Для слабого квадратичного преобразования $y=x+ax^2$ и центрированного гауссовского $x$ центрированный выход имеет
\begin{equation}
 C_y(\tau)=C_x(\tau)+2a^2C_x^2(\tau),\quad
 S_y(\omega)=S_x(\omega)+\frac{a^2}{\pi}\int S_x(\Omega)S_x(\omega-\Omega)\dd\Omega.
\end{equation}
Нечётные смешанные моменты исчезают; среднее $a\langle x^2\rangle$ вычтено. \paragraph{Спектр при распространении в недиспергирующей среде.}
Для каждой реализации гладкой простой волны в сопутствующем времени $\tau=t-z/c_0$ используем модель из~\Qref{2.8.1}:
\begin{equation}
 \partial_z u=\frac{\beta}{c_0^2}u\partial_\tau u
 =\frac{\beta}{2c_0^2}\partial_\tau(u^2).
\end{equation}
При преобразовании Фурье по $\tau$ производная даёт $-\ii\omega$, произведение --- свёртку. Следовательно,
\begin{equation}\label{eq:randomspectralevolution}
 \partial_zU(\omega,z)=-\frac{\ii\beta\omega}{4\pi c_0^2}
 \int U(\Omega,z)U(\omega-\Omega,z)\dd\Omega.
\end{equation}
Разные частоты обмениваются энергией по мере движения; это не одно статическое преобразование в точке. Для стационарного шума применяют конечное окно или спектральное представление. Дифференцируя $\langle|U|^2\rangle$, получают тройные корреляции: одним спектром задача не замыкается, а входная гауссовость не обязана сохраняться. До разрыва гладкая периодическая волна сохраняет среднее $u$ и интеграл $u^2$ (интегрирование полной производной по периоду); после разрыва нужна регуляризация Бюргерса, и возникают потери.
% Исходная модель: Цысарь, часть 2, PDF-с.10–11, формула (14).
% Преобразование и утверждение о незамкнутости моментов выведены здесь.
